\subsection{Representation: Inputs, Targets and the Unit of Prediction}
\label{sec:representation}

The representation is screened first, one factor at a time from the representation of Section~\ref{sec:algorithms}: the inputs $(\log_{10}\rhoc, \beta, \lambda)$ for neutron stars and $(\rh, \beta)$ for black holes, the targets $M$ and $Y = \log_{10}(\Dch/M)$, and one row at a time as the unit of prediction. Each alternative changes one factor: $\rhoc$ on its raw scale, the mass as $\log_{10} M$, the charge as $\Dch$ itself, or a whole curve as the unit. In the curve-wise unit a curve is fitted once as a cubic B-spline along its normalized coordinate, and a base model learns the spline coefficients and the curve's range from the curve parameters alone; a row is predicted by rebuilding its curve. Every alternative is scored at equal model through the harness of Section~\ref{sec:optimization} with three models of different kinds: a $k$-nearest-neighbour regressor weighted by inverse distance, a thin-plate radial basis function interpolant on each query's nearest neighbours, and the baseline network on three seeds. The network fits one output, so it is no curve-wise base. \Cref{tab:representation} reports, per alternative and model, the significant figures kept at the 95th percentile of the error on the test curves and on the folds; every target form is scored in $M$ or $\Dch$, so the columns compare.

\InputIfFileExists{61_representation_tab_representation}{}{}

The inputs and the targets keep the baseline's form. Raw $\rhoc$ loses between $0.14$ and $0.20$ figures on the neutron-star mass for every model, so the density enters as $\log_{10}\rhoc$, as Section~\ref{sec:eda-ns} chose from the shape of the curves. The mass as $\log_{10} M$ gains for no model on either dataset. The linear charge gains $0.19$ figures for the interpolant on neutron stars but drives the network below zero figures, a 95th-percentile error above $100\%$, while the log target keeps it at $0.6$; the charge therefore stays $Y = \log_{10}(\Dch/M)$, the form that serves every model ($H_1$ on the whole charge range; the split per decade of $\Dch$ follows with the families).

The unit of prediction does not carry across models. A whole curve as the unit gains beyond the fold and seed spread for the interpolant on the neutron-star mass, the neutron-star charge and the black-hole charge, and for the nearest neighbours on the neutron-star mass; it loses for the nearest neighbours on both black-hole targets, where at most $21$ curves inform the base. The unit is therefore not fixed here but carried into the family screen: every family is scored one row at a time and, where its model fits several outputs, curve-wise (Section~\ref{sec:families}). The best representation found, the curve-wise interpolant, keeps $3.3$ figures on the neutron-star mass and $1.3$ on its charge, against data ceilings of about $\nsEdaCeilMFigures$ and $\nsEdaCeilDFigures$; the gap the families and their tuning have to close is about one figure for each.
